\documentclass[12pt]{ctexart}
\usepackage{amsmath, amssymb, amsthm}
\usepackage{geometry}
\geometry{a4paper, margin=2.5cm}

\title{线性代数 · 第一章 · 向量与向量空间 \\ 理论作业参考答案}
\author{}
\date{}

\begin{document}
\maketitle

% ──────────────────────────────────────────────
\section*{题目 1 · 向量空间的判定}

\textbf{(1)} $W_1 = \{(x, y) \mid y = -x\}$ 是 $\mathbb{R}^2$ 的子空间。

验证三条：
\begin{itemize}
    \item 零向量：$(0, 0)$ 满足 $0 = -0$。$\checkmark$
    \item 加法封闭：设 $(x_1, -x_1), (x_2, -x_2) \in W_1$，则和 $(x_1+x_2, -(x_1+x_2)) \in W_1$。$\checkmark$
    \item 标量乘法封闭：$c(x, -x) = (cx, -cx) \in W_1$。$\checkmark$
\end{itemize}

基：$\{[1, -1]^T\}$。维数：$\dim W_1 = 1$。

\textbf{(2)} $W_2 = \{(x, y) \mid x \geq 0\}$ 不是子空间。

取 $\mathbf{v} = [1, 0]^T \in W_2$，取 $c = -1$，则 $c\mathbf{v} = [-1, 0]^T \notin W_2$（因为 $-1 < 0$）。违反了标量乘法封闭性。

\textbf{(3)} $W_3 = \{(x, y) \mid y = x^2\}$ 不是子空间。

取 $\mathbf{u} = [1, 1]^T$（$1 = 1^2$）和 $\mathbf{v} = [2, 4]^T$（$4 = 2^2$），和 $\mathbf{u} + \mathbf{v} = [3, 5]^T$。检查 $5 \neq 3^2 = 9$，所以 $\mathbf{u} + \mathbf{v} \notin W_3$。违反加法封闭性。此外 $(0,0) \in W_3$ 但仅此一条不够。

\textbf{关键教训}：只有经过原点的「直线/平面/超平面」才可能是子空间。弯曲的集合一般不封闭。

% ──────────────────────────────────────────────
\section*{题目 2 · 线性无关判定}

解方程 $c_1\mathbf{v}_1 + c_2\mathbf{v}_2 + c_3\mathbf{v}_3 = \mathbf{0}$：

\[
c_1\begin{bmatrix}1\\1\\0\end{bmatrix} +
c_2\begin{bmatrix}0\\1\\1\end{bmatrix} +
c_3\begin{bmatrix}1\\0\\-1\end{bmatrix} =
\begin{bmatrix}0\\0\\0\end{bmatrix}
\]

得方程组：
\[
\begin{cases}
c_1 + c_3 = 0 & \text{(1)} \\
c_1 + c_2 = 0 & \text{(2)} \\
c_2 - c_3 = 0 & \text{(3)}
\end{cases}
\]

由 (1)：$c_3 = -c_1$。由 (2)：$c_2 = -c_1$。代入 (3)：$-c_1 - (-c_1) = 0$ ✓。

取 $c_1 = 1$，得 $c_2 = -1, c_3 = -1$。存在非零解 $(1, -1, -1)$。

\textbf{因此向量组线性相关}。事实上 $\mathbf{v}_3 = \mathbf{v}_1 - \mathbf{v}_2$。

检查：$\mathbf{v}_1 - \mathbf{v}_2 = [1-0, 1-1, 0-1]^T = [1, 0, -1]^T = \mathbf{v}_3$ ✓。

% ──────────────────────────────────────────────
\section*{题目 3 · 求基与维数}

条件 $x_1 + x_2 + x_3 + x_4 = 0 \implies x_4 = -x_1 - x_2 - x_3$。

任意 $\mathbf{w} \in W$ 可写为：

\[
\mathbf{w} = \begin{bmatrix} x_1 \\ x_2 \\ x_3 \\ -x_1-x_2-x_3 \end{bmatrix}
= x_1\begin{bmatrix} 1 \\ 0 \\ 0 \\ -1 \end{bmatrix}
+ x_2\begin{bmatrix} 0 \\ 1 \\ 0 \\ -1 \end{bmatrix}
+ x_3\begin{bmatrix} 0 \\ 0 \\ 1 \\ -1 \end{bmatrix}
\]

候选基：$\{\mathbf{b}_1, \mathbf{b}_2, \mathbf{b}_3\} = \{[1,0,0,-1]^T, [0,1,0,-1]^T, [0,0,1,-1]^T\}$。

验证线性无关：$c_1\mathbf{b}_1 + c_2\mathbf{b}_2 + c_3\mathbf{b}_3 = [c_1, c_2, c_3, -c_1-c_2-c_3]^T = \mathbf{0} \implies c_1=c_2=c_3=0$。

故 $\dim W = 3$。

\textbf{注意}：$\mathbb{R}^4$ 中的一个方程（一个约束）使维数减少 1：$\dim W = 4 - 1 = 3$。这是「秩-零化度定理」的特例。

% ──────────────────────────────────────────────
\section*{题目 4 · 线性组合与 Span}

解 $c_1\mathbf{v}_1 + c_2\mathbf{v}_2 + c_3\mathbf{v}_3 = \mathbf{w}$：

\[
\begin{cases}
c_1 + 2c_2 + 3c_3 = 1 & \text{(1)}\\
2c_1 + 5c_2 + 3c_3 = 1 & \text{(2)}\\
c_1 + 0c_2 + 8c_3 = 1   & \text{(3)}
\end{cases}
\]

由 (3)：$c_1 = 1 - 8c_3$。

代入 (1)：$(1-8c_3) + 2c_2 + 3c_3 = 1 \implies 2c_2 - 5c_3 = 0 \implies c_2 = \frac{5}{2}c_3$。

代入 (2)：$2(1-8c_3) + 5 \cdot \frac{5}{2}c_3 + 3c_3 = 1 \implies 2 - 16c_3 + \frac{25}{2}c_3 + 3c_3 = 1$

$\implies 2 - 16c_3 + 12.5c_3 + 3c_3 = 1 \implies 2 - 0.5c_3 = 1 \implies c_3 = 2$。

回代：$c_2 = 5$, $c_1 = 1 - 16 = -15$。

验证：$-15 \cdot [1,2,1]^T + 5 \cdot [2,5,0]^T + 2 \cdot [3,3,8]^T = [-15+10+6, -30+25+6, -15+0+16]^T = [1,1,1]^T$ ✓

\textbf{所以 $\mathbf{w} = -15\mathbf{v}_1 + 5\mathbf{v}_2 + 2\mathbf{v}_3$。}

% ──────────────────────────────────────────────
\section*{题目 5 · 余弦相似度}

\textbf{(1)} 计算：
\begin{align*}
\mathbf{u} \cdot \mathbf{v} &= 1 \cdot 4 + 2 \cdot (-2) = 4 - 4 = 0 \\
\|\mathbf{u}\| &= \sqrt{1^2 + 2^2} = \sqrt{5} \approx 2.236 \\
\|\mathbf{v}\| &= \sqrt{4^2 + (-2)^2} = \sqrt{20} \approx 4.472 \\
\cos\theta &= \frac{0}{\sqrt{5} \cdot \sqrt{20}} = 0 \\
\theta &= 90.0^\circ
\end{align*}

\textbf{(2)} 余弦相似度只关心\textbf{方向}，不关心\textbf{大小}。两个文本向量即使长度差异很大（一篇长文章 vs 一条短推文），只要主题方向一致，余弦相似度就高。欧氏距离会被向量长度严重影响。

% ──────────────────────────────────────────────
\section*{题目 6 · 子空间交}

$U \cap W$ 中的向量必须同时属于两个 span，即存在标量 $\alpha_1, \alpha_2, \beta_1, \beta_2$ 使：

\[
\alpha_1\begin{bmatrix}1\\1\\0\end{bmatrix} + \alpha_2\begin{bmatrix}0\\1\\1\end{bmatrix}
= \beta_1\begin{bmatrix}1\\0\\1\end{bmatrix} + \beta_2\begin{bmatrix}0\\1\\-1\end{bmatrix}
\]

展开：
\[
\begin{bmatrix} \alpha_1 \\ \alpha_1 + \alpha_2 \\ \alpha_2 \end{bmatrix}
= \begin{bmatrix} \beta_1 \\ \beta_2 \\ \beta_1 - \beta_2 \end{bmatrix}
\]

得方程组：
\[
\begin{cases}
\alpha_1 = \beta_1 \\
\alpha_1 + \alpha_2 = \beta_2 \\
\alpha_2 = \beta_1 - \beta_2
\end{cases}
\]

代入 $\beta_1 = \alpha_1$，从第三个：$\beta_2 = \alpha_1 - \alpha_2$。代入第二个：$\alpha_1 + \alpha_2 = \alpha_1 - \alpha_2 \implies 2\alpha_2 = 0 \implies \alpha_2 = 0$。

则 $\beta_2 = \alpha_1$。向量形式为 $\alpha_1[1, 1, 0]^T = \alpha_1[1, 1, 0]^T$（两边一致）。

所以交集中的向量都是 $[1, 1, 0]^T$ 的标量倍：

\[
U \cap W = \operatorname{span}\{[1, 1, 0]^T\}
\]

基：$\{[1, 1, 0]^T\}$。维数：$\dim(U \cap W) = 1$。

% ──────────────────────────────────────────────
\section*{题目 7 · 范数不等式证明}

\textbf{(1) Cauchy-Schwarz 不等式}

考虑函数 $f(t) = \|\mathbf{u} - t\mathbf{v}\|^2 \geq 0$ 对所有 $t \in \mathbb{R}$ 成立：

\begin{align*}
f(t) &= (\mathbf{u} - t\mathbf{v}) \cdot (\mathbf{u} - t\mathbf{v}) \\
     &= \mathbf{u}\cdot\mathbf{u} - 2t(\mathbf{u}\cdot\mathbf{v}) + t^2(\mathbf{v}\cdot\mathbf{v}) \\
     &= \|\mathbf{u}\|^2 - 2t(\mathbf{u}\cdot\mathbf{v}) + t^2\|\mathbf{v}\|^2
\end{align*}

这是关于 $t$ 的二次函数（$\|\mathbf{v}\|^2 \geq 0$）。当 $\|\mathbf{v}\| \neq 0$ 时：
$f(t) \geq 0$ 对所有 $t$ 成立 $\iff$ 判别式 $\Delta \leq 0$：

\[
\Delta = 4(\mathbf{u}\cdot\mathbf{v})^2 - 4\|\mathbf{u}\|^2\|\mathbf{v}\|^2 \leq 0
\implies (\mathbf{u}\cdot\mathbf{v})^2 \leq \|\mathbf{u}\|^2\|\mathbf{v}\|^2
\implies |\mathbf{u}\cdot\mathbf{v}| \leq \|\mathbf{u}\|\|\mathbf{v}\|
\]

当 $\|\mathbf{v}\| = 0$（即 $\mathbf{v} = \mathbf{0}$），不等式两端均为 0，平凡成立。$\square$

\textbf{(2) 三角不等式}

\begin{align*}
\|\mathbf{u} + \mathbf{v}\|^2 &= (\mathbf{u} + \mathbf{v}) \cdot (\mathbf{u} + \mathbf{v}) \\
                              &= \|\mathbf{u}\|^2 + 2(\mathbf{u}\cdot\mathbf{v}) + \|\mathbf{v}\|^2 \\
                              &\leq \|\mathbf{u}\|^2 + 2\|\mathbf{u}\|\|\mathbf{v}\| + \|\mathbf{v}\|^2 \quad\text{(由 Cauchy-Schwarz)}\\
                              &= (\|\mathbf{u}\| + \|\mathbf{v}\|)^2
\end{align*}

两边开平方（范数非负），得 $\|\mathbf{u} + \mathbf{v}\| \leq \|\mathbf{u}\| + \|\mathbf{v}\|$。$\square$

% ──────────────────────────────────────────────
\section*{题目 8 · 线性无关与基的扩张}

\textbf{(1)} 已知 $\mathbf{v}_1, \mathbf{v}_2$ 线性无关。设 $c_1\mathbf{v}_1 + c_2\mathbf{v}_2 + c_3\mathbf{v}_3 = \mathbf{0}$。

若 $c_3 \neq 0$，则 $\mathbf{v}_3 = -\frac{c_1}{c_3}\mathbf{v}_1 - \frac{c_2}{c_3}\mathbf{v}_2 \in \operatorname{span}\{\mathbf{v}_1, \mathbf{v}_2\}$，与题设 $\mathbf{v}_3 \notin \operatorname{span}\{\mathbf{v}_1, \mathbf{v}_2\}$ 矛盾。故 $c_3 = 0$。

方程化为 $c_1\mathbf{v}_1 + c_2\mathbf{v}_2 = \mathbf{0}$。由 $\mathbf{v}_1, \mathbf{v}_2$ 线性无关得 $c_1 = c_2 = 0$。所以 $c_1 = c_2 = c_3 = 0$，即 $\{\mathbf{v}_1, \mathbf{v}_2, \mathbf{v}_3\}$ 线性无关。三个线性无关向量张成 $\mathbb{R}^3$（因 $\dim\mathbb{R}^3 = 3$），故是一组基。

\textbf{(2)} 需要 $\mathbf{v}_3 \notin \operatorname{span}\{\mathbf{v}_1, \mathbf{v}_2\}$。找不与 $\mathbf{v}_1, \mathbf{v}_2$ 共面的向量。

$\operatorname{span}\{\mathbf{v}_1, \mathbf{v}_2\}$ 中向量的形式：

\[
c_1\begin{bmatrix}1\\0\\-1\end{bmatrix} + c_2\begin{bmatrix}2\\1\\0\end{bmatrix}
= \begin{bmatrix} c_1+2c_2 \\ c_2 \\ -c_1 \end{bmatrix}
\]

取 $\mathbf{v}_3 = [0, 0, 1]^T$。检验它是否在 span 中：若 $[0, 0, 1]^T = [c_1+2c_2, c_2, -c_1]^T$，则 $c_2 = 0$, $-c_1 = 1 \implies c_1 = -1$, 但 $c_1+2c_2 = -1 \neq 0$，矛盾。故 $\mathbf{v}_3 \notin \operatorname{span}$。

由 (1) 知 $\{[1,0,-1]^T, [2,1,0]^T, [0,0,1]^T\}$ 是 $\mathbb{R}^3$ 的一组基。

\textbf{（也可以选 $\mathbf{v}_3 = \mathbf{v}_1 \times \mathbf{v}_2$，即叉积，天然垂直于前两个向量张成的平面，一定线性无关。）}

\end{document}
